\subsection{Macaulay 环上的局部上同调}

在本节中，我们假设 A 是局部 Macaulay 环；令 \(\dim(A) = d\) 。

由 \(\mathfrak{m}_{\mathrm{A}}\) 中对 A 完全割、且长度为 \(d = \dim (\mathrm{A})\) 的元素序列所生成的诸理想，构成子集 \(\mathcal{D}_{cs}\) ，它在 A 的定义理想组成的集合 \(\mathcal{D}\) 中共尾。事实上，设 \((x_{1},\ldots ,x_{d})\) 为 \(\mathfrak{m}_{\mathrm{A}}\) 中对 A 完全割的序列（§ 2, 第 3 节，命题 3）；对每个整数 \(n\) ，序列 \((x_1^n,\dots ,x_d^n)\) 对 A 是完全割的（A, X, 第 158 页，命题 6, c)），且生成一个定义理想（VIII, § 3, 第 2 节，命题 3 的推论及定理 1），该理想含于 \(\mathfrak{m}_{\mathrm{A}}^n\) 。

设 \(\mathfrak{a} \in \mathcal{D}_{cs}\) ，设 \(\pi: L \to A/\mathfrak{a}\) 为一个有限型的自由分解，它在次数 \(>d\) 处为零（例如与生成 \(\mathfrak{a}\) 的、对 A 完全割的序列相伴随的 Koszul 复形）。考虑 \(L^{*} = \operatorname{Homgr}_{A}(L,A)\) ，即 \(L\) 的对偶；由于 A 的深度等于 \(d\) ，有 \(\operatorname{Ext}_{A}^{i}(A/\mathfrak{a}, A) = 0\) （对 \(i < d\) ）（ \(\S\) 1, \(n^{\circ}\) 1, 命题 2 的推论 2）。由于 \(L^{*}\) 的长度 \(\leqslant d\) ，我们导出：\(H^{i}(L^{*})\) 在 \(i \neq d\) 时为零，且 \(H^{d}(L^{*})\) 等同于 \(\operatorname{Ext}_{A}^{d}(A/\mathfrak{a}, A)\) （A, X, 第 100 页，定理 1）。于是，我们有一个同态

\[
\pi^ {*}: \mathrm{L} ^ {*} (- d) \longrightarrow \operatorname{Ext} _ {\mathrm{A}} ^ {d} (\mathrm{A} / \mathfrak {a}, \mathrm{A})
\]

它定义了 \(\operatorname{Ext}_{\mathrm{A}}^{d}(\mathrm{A}/\mathfrak{a},\mathrm{A})\) 的一个有限型自由分解。

设 M 为 A-模；考虑典范同构（同前引）

\[
\varphi (\mathrm{L}, \mathrm{M}): \mathrm{H} (\operatorname{Homgr} _ {\mathrm{A}} (\mathrm{L}, \mathrm{M})) \rightarrow \operatorname{Ext} _ {\mathrm{A}} (\mathrm{A} / \mathfrak {a}, \mathrm{M})
\]

\[
\psi (\mathrm{M}, \mathrm{L} ^ {*} (- d)): \mathrm{Tor} ^ {\mathrm{A}} (\mathrm{M}, \mathrm{Ext} _ {\mathrm{A}} ^ {d} (\mathrm{A} / \mathfrak {a}, \mathrm{A})) \to \mathrm{H} (\mathrm{M} \otimes_ {\mathrm{A}} \mathrm{L} ^ {*}) (- d).
\]

由于复形 L 是有限型自由的，复形的典范态射 \(M \otimes_{A} L^{*} \to \text{Homgr}_{A}(L, M)\) 是一个同构；由此导出分次 A-模的一个同构 \(H(M \otimes_{A} L^{*}) \to H(\text{Homgr}_{A}(L, M))\) 。合成上述诸同构，我们得到分次 A-模的一个同构，它称为典范同构。

\[
\tau (\mathrm{L}, \mathrm{M}): \operatorname{Tor} ^ {\mathrm{A}} (\mathrm{M}, \operatorname{Ext} _ {\mathrm{A}} ^ {d} (\mathrm{A} / \mathfrak {a}, \mathrm{A})) (d) \longrightarrow \operatorname{Ext} _ {\mathrm{A}} (\mathrm{A} / \mathfrak {a}, \mathrm{M})
\]

它对每个整数 i 诱导一个同构

\[
\tau^ {i} (\mathrm{L}, \mathrm{M}): \operatorname{Tor} _ {d - i} ^ {\mathrm{A}} (\mathrm{M}, \operatorname{Ext} _ {\mathrm{A}} ^ {d} (\mathrm{A} / \mathfrak {a}, \mathrm{A})) \longrightarrow \operatorname{Ext} _ {\mathrm{A}} ^ {i} (\mathrm{A} / \mathfrak {a}, \mathrm{M}).
\]

对 \(M = A\) ，\(\tau^d(L, A)\) 是从 \(A \otimes_A \operatorname{Ext}_A^d(\mathrm{A} / \mathfrak{a}, A)\) 到 \(\operatorname{Ext}_\mathrm{A}^d(A / \mathfrak{a}, A)\) 上的典范同构。

设 \(\mathfrak{b}\) 为 \(\mathcal{D}_{cs}\) 中含于 \(\mathfrak{a}\) 的一个理想。设 \(\rho: R \to A/\mathfrak{b}\) 为长度 \(\leqslant d\) 的一个有限自由分解，并设 \(p_{\mathfrak{ab}}: A/\mathfrak{b} \to A/\mathfrak{a}\) 为典范满射。由 A, X, 第 49 页，命题 3，存在复形的一个态射 \(\mathrm{P}_{\mathrm{LR}}: R \to L\) ，使得 \(\pi \circ \mathrm{P}_{\mathrm{LR}} = p_{\mathfrak{ab}} \circ \rho\) 。由 A, X, 第 103 页，命题 2，我们有一个交换图

\[
\begin{array}{c c c} \operatorname{Tor} ^ {A} (M, \operatorname{Ext} _ {A} ^ {d} (A / \mathfrak {a}, A)) (d) & \xrightarrow {\operatorname{Tor} (1 _ {M} , \operatorname{Ext} ^ {d} (p _ {a b} , 1 _ {\mathrm{A}}))} & \operatorname{Tor} ^ {\mathrm{A}} (M, \operatorname{Ext} _ {A} ^ {d} (A / \mathfrak {b}, A)) (d) \\ \psi (L ^ {*} (- d), M) \Bigg \downarrow & & \Bigg \downarrow \psi (R ^ {*} (- d), M) \\ H (M \otimes_ {A} L ^ {*}) & \xrightarrow {H (1 _ {M} \otimes^ {t} P _ {L R})} & H (M \otimes_ {A} R ^ {*}) \\ \Bigg \downarrow & & \Bigg \downarrow \\ H (\operatorname{Homgr} _ {\mathrm{A}} (L, M)) & \xrightarrow {H (\operatorname{Homgr} (P _ {L R} , M))} & H (\operatorname{Homgr} _ {A} (R, M)) \\ \varphi (L, M) \Bigg \downarrow & & \Bigg \downarrow \varphi (R, M) \\ \operatorname{Ext} _ {\mathrm{A}} (A / \mathfrak {a}, M) & \xrightarrow {\operatorname{Ext} (p _ {a b} , 1 _ {M})} & \operatorname{Ext} _ {A} (A / \mathfrak {b}, M) \end{array}
\]

首先，取 \(\mathfrak{a} = \mathfrak{b}\) ，即可看出同构 \(\tau(\mathrm{L},\mathrm{M})\) 不依赖于 A/ \(\mathfrak{a}\) 的分解 L 的选取；我们把它记作 \(\tau_{\mathfrak{a}}(\mathrm{M})\) 。由此进而推出，诸 \(\tau_{\mathfrak{a}}(\mathrm{M})\) （其中 \(\mathfrak{a} \in \mathcal{D}_{cs}\) ）构成同构的一个归纳系统。取归纳极限，并注意到 A, X, 第 70 页，命题 8，对每个整数 \(i\) 得到 A-模的一个同构

\[
\tau^ {i} (\mathrm{M}): \operatorname{Tor} _ {d - i} ^ {\mathrm{A}} (\mathrm{M}, \mathrm{H} _ {\mathrm{A}} ^ {d} (\mathrm{A})) \longrightarrow \mathrm{H} _ {\mathrm{A}} ^ {i} (\mathrm{M}).
\]

对 M = A，\(\tau^{d}(A)\) 是从 \(A \otimes_{A} H_{A}^{d}(A)\) 到 \(H_{A}^{d}(A)\) 上的典范同构。

注.—— 1) 设 \(f: \mathrm{M} \to \mathrm{N}\) 为 A-模的一个同态。利用 A, X, 第 103 页，命题 2，可以证明下列诸图是交换的：

\[
\begin{array}{c c c} \operatorname{Tor} _ {d - i} ^ {\mathrm{A}} (\mathrm{M}, \mathrm{H} _ {\mathrm{A}} ^ {d} (\mathrm{A})) & \xrightarrow {\tau^ {i} (\mathrm{M})} & \mathrm{H} _ {\mathrm{A}} ^ {i} (\mathrm{M}) \\ \operatorname{Tor} _ {d - i} (f, 1) \Bigg \downarrow & & \Bigg \downarrow \mathrm{H} ^ {i} (f) \\ \operatorname{Tor} _ {d - i} ^ {\mathrm{A}} (\mathrm{N}, \mathrm{H} _ {\mathrm{A}} ^ {d} (\mathrm{A})) & \xrightarrow {\tau^ {i} (\mathrm{N})} & \mathrm{H} _ {\mathrm{A}} ^ {i} (\mathrm{N}) \quad . \end{array}
\]

\begin{enumerate}
\def\labelenumi{\arabic{enumi})}
\setcounter{enumi}{1}
\tightlist
\item
  设
\end{enumerate}

\[
0 \to \mathrm{M} \to \mathrm{N} \to \mathrm{P} \to 0\tag{ℰ}
\]

为 A-模的一个正合列。利用 A, X, 第 104 页，命题 3 以及第 106 页，命题 4，可以证明下列诸图是交换的：

\[
\begin{array}{c c c} \operatorname{Tor} _ {d - i} ^ {\mathrm{A}} (\mathrm{P}, \mathrm{H} _ {\mathrm{A}} ^ {d} (\mathrm{A})) & \xrightarrow {\tau^ {i} (\mathrm{P})} & \mathrm{H} _ {\mathrm{A}} ^ {i} (\mathrm{P}) \\ \partial_ {d - i} (\mathscr {E}, \mathrm{H} _ {\mathrm{A}} ^ {d} (\mathrm{A})) \Bigg \downarrow & & \Bigg \downarrow \partial^ {i} (\mathscr {E}) \\ \operatorname{Tor} _ {d - i - 1} ^ {\mathrm{A}} (\mathrm{M}, \mathrm{H} _ {\mathrm{A}} ^ {d} (\mathrm{A})) & \xrightarrow {\tau^ {i + 1} (\mathrm{M})} & \mathrm{H} _ {\mathrm{A}} ^ {i + 1} (\mathrm{M}) \end{array} .
\]

\begin{enumerate}
\def\labelenumi{\arabic{enumi})}
\setcounter{enumi}{2}
\tightlist
\item
  考虑同构
\end{enumerate}

\[
\tau^ {d} (\mathrm{M}): \mathrm{M} \otimes_ {\mathrm{A}} \mathrm{H} _ {\mathrm{A}} ^ {d} (\mathrm{A}) \longrightarrow \mathrm{H} _ {\mathrm{A}} ^ {d} (\mathrm{M}).
\]

对 \(x \in \mathrm{M}\) ，用 \(f_x\) 表示从 A 到 M 的映射 \(a \mapsto ax\) 。对每个 \(u \in \mathrm{H}_{\mathrm{A}}^{d}(\mathrm{A})\) ，有

\[
\tau^ {d} (\mathrm{M}) (x \otimes u) = \mathrm{H} _ {\mathrm{A}} ^ {d} (f _ {x}) (u).
\]

这确实由把注 1 应用于同态 \(f_{x}:A\to M\) 而得。

